\begin{lemma}[Existence of a finite Tashkinov closure]
Every finite tree-sequence $T$ has a closure under any fixed partial
coloring $\varphi$.
\end{lemma}
\begin{proof}
Consider all tree-sequences reachable from $T$ by a Tashkinov augmentation
sequence as defined in \noderef{TashkinovAugmentationSequence}.  This class
is nonempty because the empty augmentation sequence reaches $T$.  Every
reachable sequence has pairwise distinct vertices by \noderef{TreeSequence},
so its length is less than $|V|$.  Hence the finite set of possible lengths
has a largest member.  Choose a reachable sequence $S$ of largest length.

We claim that $V(S)$ is closed.  If it were not closed in the sense of
\noderef{ColoredVertexSetClosed}, there would be a boundary edge identity
$f$, a vertex $z\in V(S)$, and a color $c$ such that
$\varphi(f)=c$ and $c$ is missing at $z$.  By
\noderef{BoundaryEdges}, the two ends of $f$ consist of a vertex
$x\in V(S)$ and a vertex $y\notin V(S)$.  Append $f,y$ to the displayed
sequence for $S$.  The new vertex is distinct from all old vertices.  The
edge identity $f$ is distinct from every old tree edge: every old tree edge
has both ends among the old tree vertices by the attachment clause in
\noderef{TreeSequence}, whereas $f$ has the end $y$ outside that set.  The
new edge attaches $y$ to the earlier vertex $x$, and its color is missing at
the earlier vertex $z$.  The result is therefore a tree-sequence $S^+$ and
is one Tashkinov augmentation of $S$ in the sense of
\noderef{TashkinovAugmentation}.  Concatenating this step with the sequence
from $T$ to $S$ reaches $S^+$ from $T$, but $S^+$ has greater length than
$S$, contradicting maximality.

Thus $V(S)$ is closed.  Reachability together with this closure condition is
exactly \noderef{TashkinovClosure}, so $S$ is the required closure.
\end{proof}
